%--------------------------------------------------------------------------------
%
%
%--------------------------------------------------------------------------------
\pagebreak
\unnumberedSection{PxCA - Purge from Instruction or Data Cache}

\begin{iPageDescEntry}{Syntax}
    \texttt{PICA [ RegX ] ( RegB )} \\
    \texttt{PDCA [ RegX ] ( RegB )}
\end{iPageDescEntry}

\begin{iPageDescEntry}{Format}
    The \texttt{PICA} instruction uses the following instruction format: \\

    \begin{flushleft}
    \drawInstrFormat
        {0,1,3,5,6.5,8.5,9.5,11.5,16}
        {\OpCodeGrpSYS, \OpCodePCA, 0, 2, RegB, 0, RegX, 0}
    \end{flushleft}
    
    The \texttt{PDCA} instruction uses the following instruction format. \\

    \begin{flushleft}
    \drawInstrFormat
        {0,1,3,5,6.5,8.5,9.5,11.5,16}
        {\OpCodeGrpSYS, \OpCodePCA, 0, 3, RegB, 0, RegX, 0}
    \end{flushleft}
\end{iPageDescEntry}

\begin{iPageDescEntry}{Description}
    The \texttt{PxCA} instructions purge a cache line, selected by the effective address
    in \texttt{RegB}. This can be used by privileged code for cache management. Unlike \texttt{FxCA}, a purge operation may also invalidate dirty lines without writing them back.
   
\end{iPageDescEntry}

\begin{iPageDescEntry}{Operation}
    \texttt{ purgeFromCache( addOfs32( RegB, RegX ));}
\end{iPageDescEntry}

\begin{iPageDescEntry}{Exceptions}
    None.
\end{iPageDescEntry}

\begin{iPageDescEntry}{Notes}
    Typically only allowed in supervisor mode.
\end{iPageDescEntry}
